% Template for the /intro skill: a one- to two-page motivating page for the start of a course.
% Replace every <...>. Copy the preamble and macros of the course's lectures if they exist.
\documentclass[11pt,a4paper]{article}
\usepackage[margin=2.3cm]{geometry}
\usepackage{amsmath,amssymb,bm}
\usepackage{enumitem}
\usepackage{xcolor}
\usepackage{hyperref}
\hypersetup{colorlinks,linkcolor=blue!50!black,urlcolor=blue!50!black}

\newcommand{\R}{\mathbb{R}}

\title{\textbf{<Course code> <Course title>}\\[4pt]
\large Why learn <short course title>?}
\author{<Instructor>}
\date{}

\begin{document}
\maketitle

% Opening: a question or situation the students already know. One paragraph.
<Opening paragraph.>

\section*{<First use, as a plain question>}
% The real problem, what the course gives for it in everyday words, and where students will meet it.
<Two or three short paragraphs.>

\section*{<Second use>}
<...>

% 3-5 uses in all, each matching a course outcome.

\section*{How to learn this course well}
<One short paragraph.>

\end{document}
